\documentclass[10pt,a4paper]{amsart}
\usepackage[margin=19mm]{geometry}
\usepackage{amsmath,amssymb,mathtools}
\usepackage[scr=rsfs,cal=euler]{mathalfa}
\usepackage{xfrac}
\usepackage[unicode,hidelinks,bookmarks=false]{hyperref}
\newcommand{\ie}{i.e.\ }
\newcommand{\cf}{cf.\ }
\newcommand{\resp}{respectively\ }
\newcommand{\Subsection}{\subsection}
\providecommand{\nobreakdash}{\nobreak\hbox{-}}
\setlength{\emergencystretch}{2em}
\multlinegap0pt
\usepackage{bm}
\usepackage{bbm}
\usepackage{dsfont}
\usepackage{xspace}
\usepackage{cancel}
\usepackage{mathtools}
\usepackage{amsthm}
\makeatletter
\@ifundefined{theorem}{\newtheorem{theorem}{Theorem}}{}
\makeatother
\newcommand{\Aut}{\operatorname{Aut}}
\newcommand{\T}{\mathcal{T}}
\title[Isomorphisms, Moduli, and Cohomological Dimension for Twiste]{Isomorphisms, Moduli, and Cohomological Dimension for Twisted Triangular Banach Algebras}
\author{Sara Behnamian, Fatemeh Fogh}
\date{Source: arXiv:2510.03806v1 (CC BY 4.0)}
\begin{document}
\maketitle

\noindent\emph{Source and licence.} Isomorphisms, Moduli, and Cohomological Dimension for Twisted Triangular Banach Algebras. Version arXiv:2510.03806v1 (CC BY 4.0), distributed under the Creative Commons Attribution 4.0 International licence, as recorded in the arXiv licence metadata for this exact version and shown on the abstract page https://arxiv.org/abs/2510.03806v1. Full rights notice: licenses/arxiv-2510.03806-CC-BY-4.0.txt.\par\smallskip

\noindent\textit{Editorial scope.} Counted unit: the Theorem whose source label is \texttt{thm:iso} in the article (source0.tex lines 217--243, source label \texttt{thm:iso}), delivered together with the article's own complete original proof of it (source0.tex lines 245--276, one \texttt{proof} environment of 32 lines). No mathematics is written, repaired or re-derived: statement and proof are the article's own text line for line. Required context retained uncounted: none. Every definition, display and label of those lines is retained unchanged. Omitted from this package and not redistributed: 1--216, 244--244, 277--885. Nothing inside a retained excerpt is omitted or altered: every retained body is delivered exactly as the article's lines are written, so no omission receipt of any kind accompanies this package. Citations: the retained excerpts carry no citation command at all, so no bibliography entry of the article is transcribed and no reference is redistributed. Reference adaptation: none was needed and none was made; every reference inside the delivered unit resolves inside it, so no unresolved reference or citation is produced. Printed slips kept literal: no printed word, symbol, hypothesis or number of the counted statement or of its proof was altered. Layout and class: the article's own class is \texttt{article} with packages including fontenc, inputenc, lmodern, babel, geometry, microtype, amsmath, amsthm, amssymb, amsfonts, authblk, hyperref; the wrapper is the lane's pinned \texttt{amsart} editorial wrapper at 10pt on A4, which restates the theorem-like environments the retained text needs and adds the article's own macro definitions that the retained text needs (\texttt{\textbackslash Aut}, \texttt{\textbackslash T}), with extra wrapper packages (bm, bbm, dsfont, xspace, cancel, mathtools). The delivered numbering is the unit's own. Assets: no retained span contains a figure environment, a graphics-inclusion command, a PDF-page inclusion, a table or a TikZ picture; no image file is copied into this package, the package declares no asset dependency and no image file is redistributed with it. Licence: version arXiv:2510.03806v1 of this article is distributed under the Creative Commons Attribution 4.0 International licence, as recorded in the arXiv licence metadata for this exact version and shown on the abstract page https://arxiv.org/abs/2510.03806v1. A full rights notice is delivered at licenses/arxiv-2510.03806-CC-BY-4.0.txt. Characters: every retained excerpt is pure ASCII, so the delivered engine, XeLaTeX with its default Latin Modern text font, reports no missing character.
\begin{theorem}\label{thm:iso}
Let \(\sigma=(\sigma_A,\sigma_B)\) and \(\tau=(\tau_A,\tau_B)\) be twists in \(\Aut(A)\times\Aut(B)\).
Assume \(A,B\) are unital or, more generally, possess bounded approximate identities (so that the Peirce decomposition is available in the multiplier sense).
There exists a bounded algebra isomorphism
\[
\Phi:\T_\sigma(A,B;X)\longrightarrow \T_\tau(A,B;X)
\]
if and only if there are \(\alpha\in\Aut(A)\), \(\beta\in\Aut(B)\), an \((\alpha,\beta)\)--bimodule isomorphism \(u:X\to X\), and a bounded bilinear map \(\theta:A\times B\to X\) satisfying
\begin{equation}\label{eq:diag-conj}
\tau_A=\alpha\circ \sigma_A\circ \alpha^{-1},
\qquad
\tau_B=\beta\circ \sigma_B\circ \beta^{-1},
\end{equation}
and the \(\tau\)--cocycle identity
\begin{equation}\label{eq:cocycle-tau}
\theta(aa',bb')=\tau_A(\alpha(a))\cdot \theta(a',b')+\theta(a,b)\cdot \tau_B(\beta(b'))
\qquad(a,a'\in A,\ b,b'\in B).
\end{equation}
In this situation one has
\[
\Phi\!\begin{bmatrix}a&x\\[2pt]0&b\end{bmatrix}
=
\begin{bmatrix}\alpha(a)&\ u(x)+\theta(a,b)\\[2pt]0&\beta(b)\end{bmatrix}.
\]
Conversely, any data \((\alpha,\beta,u,\theta)\) satisfying \eqref{eq:diag-conj}--\eqref{eq:cocycle-tau} determine a Banach algebra isomorphism by the displayed formula.
If \(X=0\), the same description holds up to the additional possibility of interchanging the two diagonal summands.
\end{theorem}

\begin{proof}
By the characteristic property of \(I\) explained above, any isomorphism \(\Phi\) induces a homomorphism \(\bar\Phi:Q\to Q\) with \(\bar\Phi\circ\pi=\pi\circ\Phi\).
Since \(Q\cong A\oplus B\), the map \(\bar\Phi\) restricts to automorphisms \(\alpha\in\Aut(A)\) and \(\beta\in\Aut(B)\) on the two diagonal summands, unless \(X=0\) in which case a permutation of the two factors is also possible.
Write
\[
\Phi\!\begin{bmatrix}a&0\\[2pt]0&0\end{bmatrix}
=\begin{bmatrix}\alpha(a)&\theta_1(a)\\[2pt]0&0\end{bmatrix},
\qquad
\Phi\!\begin{bmatrix}0&0\\[2pt]0&b\end{bmatrix}
=\begin{bmatrix}0&\theta_2(b)\\[2pt]0&\beta(b)\end{bmatrix},
\qquad
\Phi\!\begin{bmatrix}0&x\\[2pt]0&0\end{bmatrix}
=\begin{bmatrix}0&u(x)\\[2pt]0&0\end{bmatrix},
\]
where \(u:X\to X\) is a bounded bijection and \(\theta_1,\theta_2\) are bounded linear maps.
Multiplying the first and third displays in the two possible orders and comparing off-diagonal entries shows
\[
u(\sigma_A(a)\cdot x)=\tau_A(\alpha(a))\cdot u(x)\qquad(a\in A,\ x\in X).
\]
Similarly, multiplying the third and second displays and comparing off-diagonal entries shows
\[
u(x\cdot \sigma_B(b))=u(x)\cdot \tau_B(\beta(b))\qquad(b\in B,\ x\in X).
\]
Hence \(u\) is an \((\alpha,\beta)\)-bimodule isomorphism.
Writing in general
\[
\Phi\!\begin{bmatrix}a&x\\[2pt]0&b\end{bmatrix}
=\begin{bmatrix}\alpha(a)&u(x)+\vartheta(a,b)\\[2pt]0&\beta(b)\end{bmatrix}
\]
for a bounded bilinear \(\vartheta:A\times B\to X\), multiplicativity of \(\Phi\) forces \(\vartheta\) to satisfy \eqref{eq:cocycle-tau} with \(\theta=\vartheta\), while compatibility on the diagonal forces \eqref{eq:diag-conj}.
Conversely, given \((\alpha,\beta,u,\theta)\) satisfying \eqref{eq:diag-conj}--\eqref{eq:cocycle-tau}, the displayed formula defines a bounded bijective homomorphism whose inverse is obtained from \((\alpha^{-1},\beta^{-1},u^{-1},-\theta\circ(\alpha^{-1},\beta^{-1}))\).
\end{proof}

\end{document}
